\documentclass[11pt,a4paper]{article}
\usepackage[italian]{babel}
\usepackage{amsmath,amssymb}
\usepackage[margin=2cm]{geometry}
\usepackage{enumitem}
\usepackage{array}
\usepackage{fancyhdr}
\setlength{\parindent}{0pt}
\setlength{\parskip}{2pt}
\setlength{\abovedisplayskip}{3pt}
\setlength{\belowdisplayskip}{3pt}
\setlength{\abovedisplayshortskip}{2pt}
\setlength{\belowdisplayshortskip}{2pt}
\setlist{nosep,leftmargin=1.3em}
\pagestyle{fancy}\fancyhf{}\rfoot{\thepage}
\renewcommand{\headrulewidth}{0pt}
\begin{document}
\section*{Q1. Cosa si intende per "spazio metrico" e in che modo la "metrica euclidea su $\mathbb{R}$" vi si relaziona?}
Uno \textbf{spazio metrico} è una coppia $(X, d)$ con \texttt{X} insieme non vuoto e $d: X \times  X \to  \mathbb{R}$ (distanza) che per ogni \texttt{x, y, z} soddisfa 4 assiomi:
\begin{enumerate}
\item \textbf{Non-negatività}: $d(x,y) \ge  0$
\item \textbf{Identità}: $d(x,y) = 0 \Leftrightarrow  x = y$
\item \textbf{Simmetria}: $d(x,y) = d(y,x)$
\item \textbf{Disuguaglianza triangolare}: $d(x,z) \le  d(x,y) + d(y,z)$
\end{enumerate}
La \textbf{metrica euclidea su $\mathbb{R}$} è $d(x,y) = |x - y|$: soddisfa i 4 assiomi ed è il caso prototipo. Si generalizza a $\mathbb{R}^{N}$ con $d(x,y) = \|x - y\|_{2} = \sqrt{\sum _{i}(x_{i} - y_{i})^{2}}$. La metrica serve perché tutti i concetti di limite, continuità e ottimizzazione (es. discesa del gradiente in $\mathbb{R}^{N}$) richiedono una nozione di "vicinanza".
\newpage
\section*{Q2. Come si definisce la "continuità" per una funzione @@INL35@@, e qual è il significato di "continuità separata"?}
\texttt{f} è \textbf{continua} in $x_{0}$ se:
\[ \forall \varepsilon  > 0, \exists \delta  > 0 : \|x - x_{0}\| < \delta  \Rightarrow  |f(x) - f(x_{0})| < \varepsilon  \]
Cioè: a piccole variazioni dell'input corrispondono piccole variazioni dell'output (definizione di Cauchy generalizzata da intervalli a "palle" in $\mathbb{R}^{N}$).
La \textbf{continuità separata} richiede che \texttt{f} sia continua rispetto a una variabile alla volta, tenendo fisse le altre.
\textbf{Relazione}: la continuità (congiunta) implica quella separata, \textbf{ma non viceversa}. Controesempio: $f(x,y) = xy/(x^{2}+y^{2})$ (e $f(0,0)=0$) è separatamente continua nell'origine (vale 0 sugli assi) ma non continua (vale $1/2$ sulla retta $y=x$). È la prima prova che il calcolo in più variabili non è la semplice somma di calcoli a una variabile.
\newpage
\section*{Q3. Qual è la relazione tra "derivata direzionale", "derivata parziale" e "gradiente" nel calcolo in più variabili?}
\begin{itemize}
\item \textbf{Derivata direzionale} lungo il versore \texttt{v}: $D_{v}f(x_{0}) = lim_{h\to 0} [f(x_{0}+hv) - f(x_{0})]/h$. Misura la variazione di \texttt{f} muovendosi nella direzione \texttt{v}.
\item \textbf{Derivata parziale}: derivata direzionale lungo un versore della base canonica $e_{i}$ (deriva rispetto a $x_{i}$, fissate le altre).
\item \textbf{Gradiente}: vettore delle derivate parziali $\nabla f = (\partial f/\partial x_{1}, \dots , \partial f/\partial x_{n})$.
\end{itemize}
\textbf{Relazione} (se \texttt{f} differenziabile): $D_{v}f = \nabla f \cdot  v$. Conseguenze:
\begin{itemize}
\item il gradiente contiene tutte le derivate direzionali;
\item $\nabla f$ punta nella direzione di \textbf{massima crescita} (e $-\nabla f$ di massima decrescita), per Cauchy-Schwarz;
\item $\nabla f = 0$ indica un punto critico.
\end{itemize}
Questa proprietà è il cuore della discesa del gradiente.
\newpage
\section*{Q4. Descrivi il "metodo iterativo di discesa del gradiente" e la sua finalità.}
\textbf{Finalità}: trovare un minimo di una funzione differenziabile \texttt{f} (es. la funzione di costo di una rete) quando non c'è soluzione in forma chiusa.
\textbf{Idea}: poiché $-\nabla f$ è la direzione di massima decrescita, ci si muove a ogni passo in quella direzione.
\textbf{Algoritmo}: scelto $x_{0}$ e un \textbf{learning rate} $\eta  > 0$:
\[ x_{k+1} = x_{k} - \eta  \cdot  \nabla f(x_{k}) \]
fino a un criterio di arresto ($\|\nabla f\|$ piccola, max iterazioni).
\textbf{Learning rate}: troppo grande $\to $ oscilla/diverge; troppo piccolo $\to $ lentissimo. \textbf{Garanzia}: converge al minimo globale solo se \texttt{f} è convessa; altrimenti a un punto critico. Varianti: Batch GD, SGD (un esempio per volta), Mini-batch.
\newpage
\section*{Q5. Descrivi il metodo della discesa del gradiente. In quali casi potrebbe non convergere al minimo globale?}
Metodo: $x_{k+1} = x_{k} - \eta \cdot \nabla f(x_{k})$ (vedi Q4).
Garantisce solo la convergenza a un \textbf{punto critico}. Non raggiunge il minimo globale quando:
\begin{enumerate}
\item \textbf{Funzione non convessa con minimi locali} (tipico delle reti): si ferma nel primo minimo locale.
\item \textbf{Punti di sella} ($\nabla f=0$ ma non minimo), frequenti in alta dimensione.
\item \textbf{Plateau} con gradiente quasi nullo: avanzamento bloccato (vanishing gradient).
\item \textbf{Learning rate sbagliato}: troppo grande $\to $ diverge; troppo piccolo $\to $ non arriva nel tempo utile.
\item \textbf{Inizializzazione sfortunata}: finisce nel bacino di un minimo locale scadente.
\end{enumerate}
In sintesi: minimo globale garantito \textbf{solo se \texttt{f} è convessa}.
\newpage
\section*{Q6. Come funziona un Single Layer Perceptron (SLP)? Quali sono i suoi limiti?}
L'\textbf{SLP} ha ingressi \texttt{x}, pesi \texttt{w}, bias \texttt{b} e attivazione $\varphi $ (storicamente il gradino). Calcola:
\[ z = w \cdot  x + b,   y = \varphi (z) \]
Definisce un \textbf{iperpiano} $w\cdot x + b = 0$ e classifica i punti in base al lato in cui cadono: è un \textbf{classificatore lineare binario}. Si addestra con la regola del perceptron, che converge se i dati sono linearmente separabili.
\textbf{Limiti}:
\begin{enumerate}
\item Risolve \textbf{solo problemi linearmente separabili}: NON lo XOR (Minsky-Papert).
\item Una sola superficie di decisione lineare: niente funzioni non lineari.
\item Il gradino non è differenziabile $\to $ niente backpropagation.
\end{enumerate}
Si superano con il \textbf{MLP} (strati nascosti + attivazioni non lineari).
\newpage
\section*{Q7. Confronta le funzioni di attivazione ReLU, sigmoide e tanh in termini di proprietà matematiche e applicazioni.}
\begin{center}\renewcommand{\arraystretch}{1.25}
\begin{tabular}{|l|l|l|l|}\hline
\textbf{} & \textbf{Formula} & \textbf{Codominio} & \textbf{Note} \\ \hline
\textbf{Sigmoide} & $\sigma (z)=1/(1+e^{-z})$ & (0,1) & satura $\to $ vanishing gradient; non centrata in 0 \\ \hline
\textbf{tanh} & $(e^{z}-e^{-z})/(e^{z}+e^{-z})$ & ($-1$,1) & centrata in 0; satura comunque \\ \hline
\textbf{ReLU} & \texttt{max(0,z)} & [0,$\infty $) & no saturazione per z>0; veloce; "dying ReLU" per z<0 \\ \hline
\end{tabular}\end{center}
\begin{itemize}
\item \textbf{Sigmoide}: derivata $\sigma '=\sigma (1-\sigma ) \le  1/4$; interpretabile come probabilità. Uso: output binario, gate LSTM.
\item \textbf{tanh}: output centrato $\to $ gradienti più bilanciati della sigmoide. Uso: hidden layer RNN.
\item \textbf{ReLU}: derivata 1 per z>0 $\to $ niente vanishing gradient, induce sparsità. È lo \textbf{standard} negli hidden layer delle reti deep.
\end{itemize}
\newpage
\section*{Q8. Spiega il principio dell'algoritmo di backpropagation in un MLP con un layer nascosto.}
MLP con un layer nascosto:
\[ h = \varphi (W^{(1)}x + b^{(1)}),   \hat{y} = \psi (W^{(2)}h + b^{(2)}) \]
La \textbf{backpropagation} calcola i gradienti della loss rispetto a tutti i parametri applicando la \textbf{regola della catena} all'indietro.
\textbf{Forward}: calcola $z^{(1)}, h, z^{(2)}, \hat{y}$ e la loss \texttt{L}.
\textbf{Backward} (propaga l'errore $\delta $):
\[ \delta ^{(2)} = \partial L/\partial \hat{y} \cdot  \psi '(z^{(2)})              \partial L/\partial W^{(2)} = \delta ^{(2)} h^{T} \]
\[ \delta ^{(1)} = (W^{(2)T} \delta ^{(2)}) \odot  \varphi '(z^{(1)})      \partial L/\partial W^{(1)} = \delta ^{(1)} x^{T} \]
\textbf{Update}: $W \leftarrow  W - \eta  \cdot  \partial L/\partial W$.
Riutilizzando i calcoli intermedi, il costo di un passo è \textasciitilde{}$2\times $ il forward, indipendentemente dalla profondità: è ciò che rende addestrabili gli MLP.
\newpage
\section*{Q9. Cos'è una funzione di costo e quale ruolo svolge nel processo di apprendimento? Fai un esempio con la MSE.}
Una \textbf{funzione di costo} $L(\theta )$ misura quanto le predizioni del modello si discostano dai target sul training set. L'apprendimento si formula come $\theta * = \operatorname{argmin}_\theta  L(\theta )$ e la si minimizza con la discesa del gradiente ($\nabla _\theta  L$ indica come aggiornare i parametri).
\textbf{Esempio - MSE} (regressione):
\[ L(\theta ) = (1/n) \sum _{i} (f_\theta (x_{i}) - t_{i})^{2} \]
Proprietà: differenziabile (gradiente $2(\hat{y}-t)$, ideale per backprop); penalizza quadraticamente gli errori grandi (sensibile agli outlier); minimizzarla equivale alla massima verosimiglianza con rumore gaussiano. Per la classificazione si usa invece la cross-entropy (Q29).
\newpage
\section*{Q10. Qual è il significato del teorema di approssimazione universale per le reti neurali?}
Una rete con \textbf{un solo layer nascosto} e attivazione non polinomiale, con abbastanza neuroni, può approssimare \textbf{qualunque funzione continua} su un compatto con precisione arbitraria:
\[ F(x) = \sum _{i} \alpha _{i} \varphi (w_{i}\cdot x + b_{i}),   \operatorname{sup}_K |F(x) - f(x)| < \varepsilon  \]
\textbf{Significato}: gli MLP sono \textbf{approssimatori universali}, nessuna funzione continua è fuori portata (a differenza dell'SLP, Q6).
\textbf{Cosa NON dice}: quanti neuroni servono (può essere enorme), come trovare i pesi (è un risultato di sola esistenza), nulla sulla generalizzazione. In pratica le reti \textbf{profonde} ottengono lo stesso con molti meno neuroni.
\newpage
\section*{Q11. Quali sono i componenti fondamentali di un "percettrone" o "neurone artificiale", e a cosa servono i "pesi" e il "bias"?}
Componenti: \textbf{ingressi} \texttt{x}, \textbf{pesi} \texttt{w}, \textbf{bias} \texttt{b}, \textbf{attivazione} $\varphi $. Calcolo:
\[ z = w \cdot  x + b,   y = \varphi (z) \]
\begin{itemize}
\item \textbf{Pesi}: misurano l'importanza di ogni input ($w_{i}$ grande $\to $ input influente; $w_{i}<0$ $\to $ relazione inversa). \texttt{w} è il vettore normale all'iperpiano di decisione $w\cdot x+b=0$.
\item \textbf{Bias}: \textbf{trasla} l'iperpiano; senza, passerebbe sempre per l'origine, limitando le funzioni rappresentabili. Rende la decisione \textbf{affine} invece che lineare.
\end{itemize}
L'apprendimento consiste nel trovare \texttt{w} e \texttt{b} corretti.
\newpage
\section*{Q12. Definisci la "funzione di costo" e illustra la differenza tra "loss L1" (geometria del taxi) e "loss L2" (MSE).}
Funzione di costo: $L(\theta )$ che misura l'errore medio del modello (vedi Q9).
\begin{itemize}
\item \textbf{L2 (MSE)}: $\ell  = (\hat{y}-t)^{2}$, per vettori $\|\hat{y}-t\|_{2}^{2}$. Distanza \textbf{euclidea}.
\item \textbf{L1 (MAE)}: $\ell  = |\hat{y}-t|$, per vettori $\|\hat{y}-t\|_{1}$. Distanza \textbf{Manhattan} (taxi).
\end{itemize}
\begin{center}\renewcommand{\arraystretch}{1.25}
\begin{tabular}{|l|l|l|}\hline
\textbf{} & \textbf{L1 (MAE)} & \textbf{L2 (MSE)} \\ \hline
Crescita errore & lineare & quadratica \\ \hline
Differenziabile in 0 & no & sì \\ \hline
Gradiente & $\pm 1$ costante & $2(\hat{y}-t)$ \\ \hline
Outlier & robusta & sensibile \\ \hline
Stimatore ottimo & mediana & media \\ \hline
\end{tabular}\end{center}
L2 amplifica gli errori grandi (buon segnale di convergenza ma sensibile agli outlier); L1 è robusta ma con gradiente costante. Compromesso: \textbf{Huber loss}.
\newpage
\section*{Q13. Elenca e descrivi brevemente almeno quattro "funzioni di attivazione" comunemente utilizzate nelle reti neurali artificiali.}
L'attivazione $\varphi $ introduce non linearità (senza, la rete collasserebbe in una trasformazione lineare).
\begin{enumerate}
\item \textbf{Sigmoide} $\sigma (z)=1/(1+e^{-z})$ $\in $ (0,1): interpretabile come probabilità; satura (vanishing gradient). Uso: output binario.
\item \textbf{tanh} $\in $ ($-1$,1): come la sigmoide ma centrata in 0 $\to $ gradienti più bilanciati. Uso: hidden layer RNN.
\item \textbf{ReLU} \texttt{max(0,z)}: nessuna saturazione per z>0, calcolo banale, sparsità. Standard negli hidden layer.
\item \textbf{Softmax} $e^{zi}/\sum _{j}e^{zj}$: trasforma un vettore in distribuzione di probabilità (somma 1). Uso: output classificazione multiclasse.
\end{enumerate}
(Varianti: Leaky ReLU, GELU.)
\newpage
\section*{Q14. Qual è il principio su cui si basa l'algoritmo di "backpropagation" e come mai è importante per l'addestramento di MLP?}
Si basa su due idee:
\begin{enumerate}
\item \textbf{Regola della catena}: una rete è una composizione di funzioni, quindi il gradiente è il prodotto delle derivate locali lungo gli strati.
\item \textbf{Calcolo all'indietro}: si propaga il "segnale di errore" dall'output all'input riusando i calcoli intermedi:
\end{enumerate}
\[ \delta ^{(l)} = (W^{(l+1)T} \delta ^{(l+1)}) \odot  \varphi '(z^{(l)}) \]
da cui $\partial L/\partial W^{(l)}$ e $\partial L/\partial b^{(l)}$.
\textbf{Importanza}: calcola \textbf{tutti} i gradienti in tempo $O(P)$, come un solo forward pass (senza, servirebbero \texttt{P} valutazioni con differenze finite $\to $ training impossibile per reti con milioni di parametri). Funziona per qualunque architettura differenziabile e ha reso possibile il deep learning.
\newpage
\section*{Q15. Cosa afferma il "Teorema di approssimazione universale"?}
Afferma che una rete feedforward con \textbf{un solo layer nascosto} e attivazione non polinomiale (sigmoide, tanh, ReLU…), con un numero sufficiente di neuroni, può approssimare \textbf{qualunque funzione continua su un compatto} con precisione arbitraria:
\[ F(x) = \sum _{i} \alpha _{i} \varphi (w_{i}\cdot x + b_{i}),   \operatorname{sup}_K |F(x) - f(x)| < \varepsilon  \]
È la \textbf{giustificazione teorica} delle reti neurali come modelli generali. Garantisce l'esistenza di pesi buoni, ma \textbf{non} quanti neuroni servano, né che la discesa del gradiente li trovi, né la generalizzazione su dati nuovi. (Stesso contenuto di Q10.)
\newpage
\section*{Q16. Descrivere le fasi principali del processo di "digitalizzazione delle immagini", specificando i concetti di "campionamento" e "quantizzazione".}
La digitalizzazione converte un'immagine continua in una matrice di pixel in due fasi principali:
\begin{itemize}
\item \textbf{Campionamento}: discretizzazione del \textbf{dominio spaziale}. Si valuta l'immagine su una griglia $M\times N$ di punti (i pixel) $\to $ definisce la \textbf{risoluzione}. Campionare troppo poco causa \textbf{aliasing}; il teorema di Nyquist richiede frequenza di campionamento $\ge $ doppio della massima frequenza del segnale.
\item \textbf{Quantizzazione}: discretizzazione del \textbf{codominio} (intensità). I valori continui vengono mappati su \texttt{L} livelli (tipicamente $L=256$, 8 bit). Pochi livelli $\to $ \textbf{banding/posterizzazione}.
\end{itemize}
Risultato: immagine digitale $I \in  {0,\dots ,L-1}^{M\times N\times C}$ ($C=1$ grigi, $C=3$ RGB).
\newpage
\section*{Q17. Cosa si intende per immagine digitale? Spiega i concetti di campionamento, quantizzazione e spazio colorimetrico.}
Un'\textbf{immagine digitale} è una matrice (tensore) di pixel $H\times W\times C$, ottenuta per digitalizzazione.
\begin{itemize}
\item \textbf{Campionamento}: discretizzazione spaziale su griglia $\to $ risoluzione (vedi Q16).
\item \textbf{Quantizzazione}: discretizzazione dei valori di intensità su \texttt{L} livelli (vedi Q16).
\item \textbf{Spazio colorimetrico}: il sistema con cui i numeri del pixel diventano colore:
\item \textbf{RGB}: additivo, 3 canali, standard per i display;
\item \textbf{Scala di grigi}: 1 canale, $Y = 0.299R + 0.587G + 0.114B$ (luminanza);
\item \textbf{HSV}: tinta/saturazione/valore, utile per selezioni per colore;
\item \textbf{YCbCr}: separa luminanza e crominanza, sfruttato nella compressione (JPEG).
\end{itemize}
\newpage
\section*{Q18. Come si utilizza un istogramma per operazioni di elaborazione su immagini in scala di grigi?}
L'\textbf{istogramma} $h(k)$ conta i pixel con intensità \texttt{k}; normalizzato $p(k)=h(k)/(MN)$ è la distribuzione delle intensità. Rivela esposizione e contrasto (stretto = basso contrasto; spostato = sotto/sovra-esposto).
Operazioni:
\begin{itemize}
\item \textbf{Stretching del contrasto}: rimappa il range usato \texttt{[a,b]} su $[0,L-1]$.
\item \textbf{Equalizzazione}: usa la CDF come trasformazione, $I' = (L-1)\cdot CDF(I)$ $\to $ istogramma \textasciitilde{}uniforme, massimizza il contrasto.
\item \textbf{Sogliatura (thresholding)}: scelta una soglia \texttt{T} (anche automatica, Otsu) si binarizza per separare oggetto/sfondo.
\end{itemize}
Limite: l'istogramma \textbf{ignora la struttura spaziale} (due immagini diverse possono avere lo stesso istogramma).
\newpage
\section*{Q19. Spiega come funziona il filtro di Sobel e quale informazione estrae da un'immagine.}
Il \textbf{Sobel} approssima il \textbf{gradiente} dell'intensità per estrarre i \textbf{bordi}. Usa due kernel $3\times 3$:
\begin{verbatim}
        | -1  0  +1 |              | -1  -2  -1 |
S_x =   | -2  0  +2 |       S_y =  |  0   0   0 |
        | -1  0  +1 |              | +1  +2  +1 |
\end{verbatim}
Per convoluzione si ottengono $G_x = I*S_x$ e $G_y = I*S_y$, da cui:
\[ M = \sqrt{G_x^{2} + G_y^{2}}        (magnitudo: forza del bordo) \]
\[ \theta  = \operatorname{arctan2}(G_y, G_x)        (direzione del bordo) \]
\texttt{M} alta = bordo, bassa = zona uniforme. I pesi \texttt{(1,2,1)} fanno anche uno \textbf{smoothing} implicito $\to $ meno sensibile al rumore del Prewitt. È il secondo stadio della pipeline di Canny (progetto d'esame).
\newpage
\section*{Q20. Qual è la differenza tra convoluzione e cross-correlazione? Fornisci una definizione formale.}
\textbf{Cross-correlazione} (2D): il kernel scorre e si applica direttamente:
\[ (I \star  K)[i,j] = \sum _{u,v} I[i+u, j+v] \cdot  K[u,v] \]
\textbf{Convoluzione} (2D): come sopra ma con il \textbf{kernel ribaltato} (segno meno):
\[ (I * K)[i,j] = \sum _{u,v} I[i-u, j-v] \cdot  K[u,v] \]
Differenza: la convoluzione ribalta il kernel di $180^{\circ}$ prima di applicarlo; per kernel simmetrici coincidono. Solo la \textbf{convoluzione} è commutativa e associativa (e soddisfa il teorema di Fourier). Le CNN, nonostante il nome, usano la cross-correlazione: i pesi sono appresi, quindi il flip è irrilevante.
\newpage
\section*{Q21. Quali sono le proprietà principali dell'operazione di convoluzione?}
\begin{enumerate}
\item \textbf{Commutativa}: $f*g = g*f$
\item \textbf{Associativa}: $(f*g)*h = f*(g*h)$ $\to $ permette di comporre filtri
\item \textbf{Distributiva} sulla somma: $f*(g+h) = f*g + f*h$
\item \textbf{Lineare}: $(\alpha f)*g = \alpha (f*g)$
\item \textbf{Elemento neutro}: la delta di Dirac, $f*\delta  = f$
\item \textbf{Differenziazione}: $(f*g)' = f'*g = f*g'$ (es. derivata della gaussiana)
\item \textbf{Teorema di convoluzione}: $F{f*g} = F{f}\cdot F{g}$ $\to $ in frequenza diventa prodotto, calcolo via FFT
\item \textbf{Invarianza per traslazione} $\to $ base dell'equivarianza delle CNN
\item \textbf{Separabilità}: se $K = Kx\cdot Ky$, costo da $O(k^{2})$ a $O(k)$ (es. gaussiana)
\end{enumerate}
\newpage
\section*{Q22. Illustra l'utilizzo dell'"istogramma dei livelli di grigio" in relazione agli "operatori puntuali" per la manipolazione di immagini luminanza.}
Un \textbf{operatore puntuale} trasforma ogni pixel in funzione del solo suo valore: $I'[i,j] = T(I[i,j])$, senza guardare i vicini. L'effetto sull'istogramma è prevedibile: se \texttt{T} è iniettiva $h'(T(k)) = h(k)$.
Principali operatori (definiti tramite \texttt{T}):
\begin{itemize}
\item \textbf{Negativo}: $T(k) = (L-1) - k$
\item \textbf{Correzione gamma}: $T(k) = (L-1)(k/(L-1))^\gamma $ ($\gamma $<1 schiarisce le ombre)
\item \textbf{Contrast stretching}: $T(k) = (L-1)(k-a)/(b-a)$
\item \textbf{Equalizzazione}: $T(k) = (L-1)\cdot CDF(k)$ $\to $ istogramma uniforme
\item \textbf{Sogliatura}: $T(k) = 0$ se $k<T$, altrimenti $L-1$
\end{itemize}
Si implementano via \textbf{lookup table} (\texttt{O(1)} per pixel). Limite: ciechi al contesto spaziale (non riducono rumore, non rilevano bordi).
\newpage
\section*{Q23. In che modo i "polinomi di Taylor" vengono applicati per la creazione di "filtri di sharpening" come il filtro di Sobel o il filtro Laplaciano?}
Le derivate discrete di un'immagine si ricavano dallo \textbf{sviluppo di Taylor} dei pixel vicini.
\textbf{Derivata prima} (Sobel): sottraendo gli sviluppi di $I[i,j+1]$ e $I[i,j-1]$:
\[ \partial I/\partial x \approx  (I[i,j+1] - I[i,j-1]) / 2     \to  kernel [-1, 0, +1] \]
mediato sulle tre righe con pesi \texttt{(1,2,1)} dà il kernel di Sobel.
\textbf{Derivata seconda} (Laplaciano): sommando gli sviluppi:
\[ \partial ^{2}I/\partial x^{2} \approx  I[i,j+1] - 2I[i,j] + I[i,j-1] \]
\[ \nabla ^{2}I = \partial ^{2}I/\partial x^{2} + \partial ^{2}I/\partial y^{2}   \to  kernel [[0,1,0],[1,-4,1],[0,1,0]] \]
\textbf{Sharpening}: $I_sharp = I - k\cdot \nabla ^{2}I$ (accentua i bordi). In sintesi: i filtri differenziali sono \textbf{polinomi di Taylor discretizzati} sui vicini di un pixel.
\newpage
\section*{Q24. Definisci l'operazione di convoluzione nel continuo e indicane le principali proprietà, sia in una che in due variabili.}
\textbf{1 variabile}:
\[ (f * g)(t) = \int  f(\tau ) g(t - \tau ) d\tau  \]
\textbf{2 variabili}:
\[ (f * g)(x,y) = \iint  f(u,v) g(x-u, y-v) du dv \]
Concettualmente: si fa scorrere \texttt{g} (riflessa) su \texttt{f} e si integra il prodotto (media pesata).
Proprietà: commutativa, associativa, distributiva, lineare; elemento neutro $\delta $; differenziazione $(f*g)' = f'*g$; teorema di convoluzione (Fourier $\to $ prodotto); invarianza per traslazione; separabilità in 2D. (Vedi Q21.) È un'operazione \textbf{lineare e invariante per traslazione}, base di filtri, sistemi LTI e CNN.
\newpage
\section*{Q25. Quali sono le differenze strutturali e funzionali tra un MLP e una CNN?}
\begin{center}\renewcommand{\arraystretch}{1.25}
\begin{tabular}{|l|l|l|}\hline
\textbf{Aspetto} & \textbf{MLP} & \textbf{CNN} \\ \hline
Connettività & fully-connected & locale (kernel) \\ \hline
Pesi & indipendenti & \textbf{condivisi} \\ \hline
N. parametri & altissimo & ridotto \\ \hline
Input & vettore 1D & tensore $H\times W\times C$ \\ \hline
Equivarianza traslazione & no & sì \\ \hline
\end{tabular}\end{center}
L'MLP appiattisce l'immagine perdendo la struttura spaziale; la CNN preserva la griglia 2D ed estrae feature locali con pesi condivisi. Su un'immagine $224\times 224\times 3$, un MLP avrebbe \textasciitilde{}$10^{8}$ parametri nel primo strato, una CNN poche migliaia. Una CNN è di fatto un MLP con vincoli (località + condivisione pesi) che incorporano l'\textbf{inductive bias} delle immagini: molti meno parametri ed equivarianza per traslazione.
\newpage
\section*{Q26. Perché le "Convolutional Neural Networks (CNN)" sono particolarmente efficaci per la gestione di "dati a griglia", come le immagini?}
Sfruttano per costruzione tre proprietà dei dati a griglia:
\begin{enumerate}
\item \textbf{Località}: i pattern (bordi, texture) sono in regioni locali $\to $ kernel piccoli; strati profondi vedono regioni più grandi (campo recettivo crescente).
\item \textbf{Condivisione dei pesi}: lo stesso kernel ovunque $\to $ \textbf{equivarianza per traslazione}, pochi parametri (un kernel $3\times 3$ ha $9\cdot C\cdot C'$ pesi indipendentemente dalla risoluzione), generalizzazione spaziale.
\item \textbf{Composizionalità gerarchica}: con conv + pooling si costruiscono feature sempre più astratte (bordi $\to $ parti $\to $ oggetti).
\end{enumerate}
Imporre queste simmetrie come inductive bias (invece di farle apprendere a un MLP) riduce drasticamente i parametri e migliora la generalizzazione.
\newpage
\section*{Q27. Dare una breve descrizione dei minimi esempi di strati utilizzati in una CNN: "layer convoluzionale", "funzione di attivazione ReLU", "layer di pooling", e "dropout".}
\begin{itemize}
\item \textbf{Convoluzionale}: applica \texttt{K} filtri appresi (es. $3\times 3$) all'input $\to $ feature map. $y[i,j,k] = \sum  W_k[u,v,c]\cdot x[i+u,j+v,c] + b_k$. Estrae feature locali. Iperparametri: kernel size, stride, padding, n. filtri.
\item \textbf{ReLU}: \texttt{max(0,z)} element-wise dopo la conv. Introduce non linearità, niente vanishing gradient per z>0, induce sparsità.
\item \textbf{Pooling} (es. max $2\times 2$): riduce la risoluzione spaziale $\to $ meno calcolo, invarianza a piccole traslazioni, campo recettivo maggiore. Senza parametri.
\item \textbf{Dropout}: in training spegne ogni neurone con probabilità \texttt{p} (e riscala di $1/(1-p)$). \textbf{Regolarizzazione} anti-overfitting (effetto ensemble). Disattivato in inference.
\end{itemize}
Stack tipico: $Conv \to  ReLU \to  Pool \to  \dots  \to  FC \to  Dropout \to  FC \to  Softmax$.
\newpage
\section*{Q28. Per quale tipo di dati sono concepite le "Recurrent Neural Networks (RNN)" e quale problematica affrontano?}
Le \textbf{RNN} sono concepite per \textbf{dati sequenziali} (testo, audio, serie temporali, video) in cui l'ordine conta e ogni elemento dipende dal contesto precedente.
Risolvono due limiti delle reti feedforward sulle sequenze: input di lunghezza variabile e assenza di memoria. Introducono uno \textbf{stato nascosto} aggiornato a ogni passo:
\[ h_t = \varphi (W_h h_{t-1} + W_x x_t + b) \]
\texttt{h\_t} riassume il passato; i pesi sono \textbf{condivisi nel tempo}. Si addestrano con la backpropagation through time.
\textbf{Problema principale}: \textbf{vanishing/exploding gradient} su lunghe sequenze (il gradiente attraversa \texttt{T} moltiplicazioni per \texttt{W\_h}) $\to $ non imparano dipendenze a lungo termine. Soluzione: \textbf{LSTM/GRU}, che con dei gate mantengono selettivamente la memoria.
\newpage
\section*{Q29. Quali sono le principali caratteristiche della "funzione di costo cross-entropy" e della "funzione di attivazione softmax", e in quale contesto sono spesso utilizzate?}
Coppia standard per la \textbf{classificazione multiclasse}.
\textbf{Softmax} trasforma i logit in distribuzione di probabilità:
\[ \operatorname{softmax}(z)_{i} = e^{zi} / \sum _{j} e^{zj},    \sum _{i} = 1 \]
Generalizza la sigmoide a \texttt{K} classi; invariante per costante additiva (si sottrae $\operatorname{max}(z)$ per stabilità).
\textbf{Cross-entropy} misura la distanza dalla distribuzione vera; con target one-hot:
\[ L = - \sum _{i} t_{i} \operatorname{log} p_{i} = - \operatorname{log} p_c   (c = classe vera) \]
È la negative log-likelihood (massima verosimiglianza).
Insieme danno un \textbf{gradiente pulito}: $\partial L/\partial z_{i} = p_{i} - t_{i}$. Uso: classificazione multiclasse, modelli di linguaggio (next token). Per $K=2$ $\to $ binary cross-entropy con sigmoide.
\newpage
\section*{Q30. Descrivi l'approccio dei modelli "Generative Adversarial Networks (GAN)".}
Le \textbf{GAN} sono modelli generativi basati su due reti in \textbf{competizione}:
\begin{itemize}
\item \textbf{Generatore} \texttt{G}: da rumore \texttt{z} genera dati falsi $G(z)$, vuole ingannare D.
\item \textbf{Discriminatore} \texttt{D}: stima la probabilità che un campione sia reale.
\end{itemize}
Si addestrano in un \textbf{gioco minimax}:
\[ \operatorname{min}_G \operatorname{max}_D  E[\operatorname{log} D(x)] + E[\operatorname{log}(1 - D(G(z)))] \]
\texttt{D} massimizza (distinguere vero/falso), \texttt{G} minimizza (produrre campioni credibili). All'equilibrio \texttt{p\_G = p\_data} e $D = 1/2$ (non distingue più).
\textbf{Difficoltà}: addestramento instabile, \textbf{mode collapse} (G genera pochi campioni ripetitivi), vanishing gradient per G se D è troppo bravo (si usa la non-saturating loss $-\operatorname{log} D(G(z))$).
\newpage
\section*{Q31. Quali problemi possono sorgere nella retropropagazione in MLP profondi (vanishing/exploding gradients)?}
Nella backprop il gradiente di uno strato profondo è il \textbf{prodotto} di molte Jacobiane $J^{(k)} = W^{(k)T} \operatorname{diag}(\varphi '(z^{(k)}))$. Se la norma di questa catena cresce o decresce esponenzialmente con la profondità:
\begin{itemize}
\item \textbf{Vanishing}: $\|\partial L/\partial W\| \to  0$, gli strati iniziali non imparano. Cause: attivazioni saturanti (sigmoide/tanh, $\varphi '\approx 0$), pesi piccoli.
\item \textbf{Exploding}: $\|\partial L/\partial W\| \to  \infty $, aggiornamenti enormi, divergenza/NaN. Cause: pesi grandi, RNN su lunghe sequenze.
\end{itemize}
\textbf{Soluzioni}: inizializzazione \textbf{Xavier/He}; attivazioni non saturanti (\textbf{ReLU}); \textbf{Batch Normalization}; \textbf{skip/residual connections} (ResNet); \textbf{gradient clipping} (per l'exploding); architetture con gate (LSTM); optimizer adattivi (Adam). Il deep learning moderno è in gran parte l'insieme di queste tecniche.
\newpage
\section*{Q32. Spiega il funzionamento di una CNN, indicando il ruolo di ciascun tipo di layer (convoluzionale, pooling, lineare).}
Una \textbf{CNN} trasforma un'immagine $H\times W\times C$ in un output (es. classi) tramite una \textbf{gerarchia di feature}: bordi $\to $ parti $\to $ oggetti. Ruolo dei layer:
\begin{itemize}
\item \textbf{Convoluzionale}: estrae \textbf{feature locali} con filtri condivisi ($y = \sum W_k\cdot x + b_k$). Località + condivisione pesi $\to $ equivarianza e pochi parametri. Seguito da ReLU.
\item \textbf{Pooling}: \textbf{down-sampling} spaziale; riduce calcolo, dà invarianza a piccole traslazioni, aumenta il campo recettivo. Senza parametri.
\item \textbf{Lineare (fully-connected)}: alla fine, combina globalmente le feature per la \textbf{classificazione} (+ softmax). Sta in fondo perché ha molti parametri e lavora su feature già compatte.
\end{itemize}
Struttura tipica:
\begin{verbatim}
Input -> [Conv -> ReLU -> Pool] x N -> Flatten -> FC -> Softmax
\end{verbatim}
*Fine delle 32 domande.*
\end{document}